% Generated draft: requires mathematical and source review.
\documentclass[11pt]{article}
\usepackage[margin=25mm]{geometry}
\usepackage{fontspec}
\setmainfont{Latin Modern Roman}
\usepackage{amsmath,amssymb,mathtools,unicode-math}
\setmathfont{Latin Modern Math}
\usepackage{enumitem,needspace,xurl,hyperref}
\setlength{\parindent}{0pt}
\setlength{\parskip}{5pt}
\setlength{\emergencystretch}{3em}
\newcommand{\N}{\mathbb{N}}
\newcommand{\Z}{\mathbb{Z}}
\newcommand{\Q}{\mathbb{Q}}
\newcommand{\R}{\mathbb{R}}
\newcommand{\C}{\mathbb{C}}
\newcommand{\F}{\mathbb{F}}
\newcommand{\A}{\mathbb{A}}
\newcommand{\PP}{\mathbb{P}}
\newcommand{\E}{\mathbb{E}}
\newcommand{\eps}{\varepsilon}
\newcommand{\dd}{\,\mathrm{d}}
\begin{document}
\section*{UnsolvedMath ID: 6000011}
\textbf{Source problem code:} catalogue entry
\subsection{Definitions and mathematical statement}
In the usual convention let $(M^n,g,\nabla)$ be a statistical manifold: $g$ is Riemannian, $\nabla$ is torsion free, and $\nabla g$ is totally symmetric. The dual connection is determined by
\[
 Xg(Y,Z)=g(\nabla_XY,Z)+g(Y,\nabla_X^*Z).
\]
The dual statistical manifold is $(M,g,\nabla^*)$.

For a point $p\in M$, choose a $\nabla$-normal neighborhood on which the geodesic from $p$ to $q$ is uniquely determined. For $q\ne p$ in that neighborhood let $\gamma_{p,q}(0)=p$, $\gamma_{p,q}(1)=q$ be that affinely parametrized geodesic and put
\[
 V_p(q)=\dot\gamma_{p,q}(1),\qquad
 D_p(q)=\{w\in T_qM:g_q(w,V_p(q))=0\}.
\]
Where $V_p(q)\ne0$, this is a rank-$(n-1)$ distribution. Integrability means that its tangent subspaces are tangent to local codimension-one leaves, equivalently that local sections are closed under Lie bracket.

The supplied implication asks whether integrability of every such radial distribution forces the dual statistical manifold to be one-conformally flat. A local reading of the hypothesis is: for every center $p$ and each domain where the radial field is smoothly defined, $D_p$ is integrable. The source does not say whether this local reading or a stronger global construction across multiple geodesics is intended.

One-conformal flatness is not defined in the record. It is a specialized statistical-geometric condition, not silently replaced here by ordinary Riemannian conformal flatness or projective flatness. Its metric/connection transformation, local versus global meaning, and dimensional hypotheses must be supplied before the conclusion is fully explicit. The source also does not state a lower bound on $n$ or address the singularity of the radial construction at $p$.
\subsection{Short English statement}
If every distribution orthogonal to geodesics radiating from a point is integrable, must the dual statistical structure be one-conformally flat? The radial field is defined above on normal neighborhoods, while the specialized conclusion remains an unresolved source definition.

The supplied research summary reports no verified published proof or counterexample and places the question in the conformal/projective-flatness theme of statistical geometry. It gives no precise partial theorem despite the registry's partial-progress label. The 1998 source presents it as open; provenance marks current status for review.
\subsection{Sources}
\begin{itemize}
\item Local UnsolvedMath registry, permanent ID 6000011, source code catalogue entry, ``Integrable Radial Distributions and One-Conformal Flatness,'' proposed by Takashi Kurose, 1998. Dataset curation and metadata: Ulam AI and UnsolvedMath Contributors, CC BY 4.0.
\item Furuhata, Matsuzoe, and Urakawa, \emph{Open Problems in Affine Differential Geometry and Related Topics} (1998), item 3(e), scanned PDF pages 1--2, \url{https://www.jstage.jst.go.jp/article/iis/4/2/4_2_125/_pdf/-char/en}; manual OCR from page images, accessed 2026-07-29.
\item The supplied summary provides no later primary reference, exact progress theorem, or definition of one-conformal flatness.
\item Source-specific provenance: the article presents this as an open problem and current status NEEDS\_REVIEW. The J-STAGE article page marks the article CC BY 4.0; this attribution is retained in addition to the dataset attribution.
\end{itemize}
\end{document}
